\documentclass[10pt]{article}
\usepackage[margin=1in]{geometry}
\usepackage{amsmath,amssymb,amsthm,hyperref}
\newtheorem{theorem}{Theorem}
\newtheorem{lemma}{Lemma}
\title{MathDB 367323: all multiples of three,\ with the exact signs}
\author{Principia Math In Progress}
\date{18 September 2026}
\begin{document}
\maketitle

\section*{Task 1: Complete intended determinant formulas}
Write $M(z)=1+zM(z)+z^2M(z)^2$ and
\[
 H_N^{(r)}=\det\bigl([z^{i+j}]M(z)^r\bigr)_{0\le i,j<N},\qquad H_0^{(r)}=1.
\]
Wang and Zhang's Conjecture 3
(\href{https://arxiv.org/abs/2502.21050}{arXiv:2502.21050}, p.~2)
has a sign error: $H_9^{(9)}=256$, whereas its printed sign gives
$-256$. We prove both intended all-parameter formulas, including
their precise signs. The result is substantially more than this
isolated correction.

\begin{theorem}\label{target}
Let $r\ge3$ be divisible by three and put
$\epsilon_r=(-1)^{\binom r2}$. For every integer $n\ge0$,
\begin{align}
 H_{rn}^{(r)}=H_{rn+1}^{(r)}&=\epsilon_r^n(n+1)^{r-1},\label{main1}\\
 H_{rn-1}^{(r)}+H_{rn+2}^{(r)}
 &=-\frac{r(r-3)}2\epsilon_r^n(n+1)^{r-1}.\label{main2}
\end{align}
For $n=0$ the second identity uses the harmless convention
$H_{-1}^{(r)}=0$; without it, take $n\ge1$ there.
\end{theorem}
Thus the absolute value required of the source's $\alpha$ is
correct, with $\alpha=-\epsilon_r^n r(r-3)/18$. The sign in
\eqref{main1} is $(-1)^n$ for $r\equiv3,6\pmod{12}$ and $1$
for $r\equiv0,9\pmod{12}$.

We prove a stronger polynomial identity. For a real parameter $a$,
let $M_a(z)=1+azM_a(z)+z^2M_a(z)^2$, and denote the corresponding
determinants by $H_N^{(r)}(a)$. Define
\[
 P_j(t)=U_j\!\left(\frac{t-a}{2}\right),\qquad
 A(t)=P_{r-1}(t),\qquad
 B_k(t)=U_k\!\left(T_r\!\left(\frac{t-a}{2}\right)\right),
\]
where $T_j(\cos\theta)=\cos(j\theta)$ and
$U_j(\cos\theta)=\sin((j+1)\theta)/\sin\theta$.
The polynomials $P_j$ and $B_k$ are monic, of degrees $j$ and $rk$.

\begin{theorem}[Weighted block boundary identities]\label{weighted}
For every $r\ge2$, $k\ge0$, and real or complex $a$, with $N=rk$,
\begin{align}
 H_N^{(r)}(a)=H_{N+1}^{(r)}(a)
 &=\epsilon_r^k B_k(0)^{r-1},\label{weighted1}\\
 H_{N-1}^{(r)}(a)+H_{N+2}^{(r)}(a)
 &=\left(r-\binom r2a^2\right)H_N^{(r)}(a).\label{weighted2}
\end{align}
Again $H_{-1}^{(r)}(a)=0$.
\end{theorem}

\subsection*{The convolution functional and its determinant}
For real $a$, let $\mu_a$ be the translated semicircle measure
\[
 d\mu_a(t)=\frac{\sqrt{4-(t-a)^2}}{2\pi}\,dt\quad(a-2\le t\le a+2).
\]
Its moments are $[z^j]M_a(z)$, and its monic orthogonal polynomials
are $P_j$, with $\int P_iP_j\,d\mu_a=\delta_{ij}$. Both statements
follow by $t=a+2\cos\theta$: the centered moments are the Catalan
numbers and the orthogonality integral is the sine integral.
Set $Q(t)=t^{r-1}A(t)$ and $L(f)=\int f(t)Q(t)\,d\mu_a(t)$.
Then
\begin{equation}\label{functional}
 [z^j]M_a(z)^r=L(t^j).
\end{equation}
Indeed $G(w)=w^{-1}M_a(w^{-1})$ satisfies $G^2=(w-a)G-1$,
so induction gives $G^r=P_{r-1}(w)G-P_{r-2}(w)$. Multiplication
by $w^{r-1}$ and extraction of $w^{-j-1}$ gives \eqref{functional}.
The functional $L$ need not be positive.

We use the following classical Christoffel identity in an algebraic
form. Its evaluation form is Theorem 1 of Krattenthaler,
\href{https://arxiv.org/abs/2101.04225}{arXiv:2101.04225}, with
repeated roots covered by Proposition 5. A proof is included to make
the argument self-contained.

\begin{lemma}\label{remainder}
If $Q$ is monic of even degree $d$, the determinant of the moments
of $Q\,d\mu_a$ of order $m$ is the determinant of the $d\times d$
matrix whose $j$th column is the coefficient vector of
$P_{m+j}\bmod Q$, for $0\le j<d$.
\end{lemma}
\begin{proof}
First suppose $Q$ has distinct roots $\alpha_1,\ldots,\alpha_d$.
Expansion of determinants gives the moment determinant as
\[
 \frac1{m!}\int\cdots\int\Delta(x)^2\prod_{i=1}^m Q(x_i)
 \prod_{i=1}^m d\mu_a(x_i).
\]
The mixed Vandermonde satisfies
$\Delta(x,\alpha)=\Delta(x)\Delta(\alpha)\prod_{i,j}(\alpha_j-x_i)$.
Replace monomials in both Vandermondes by the monic $P_j$.
Orthogonality leaves only columns $0,\ldots,m-1$ for the $x$ rows,
with integral $m!$. Hence the determinant is
\[
 (-1)^{md}\frac{\det(P_{m+j}(\alpha_i))_{1\le i\le d,\,0\le j<d}}
 {\det(\alpha_i^j)_{1\le i\le d,\,0\le j<d}}.
\]
The sign is one. Evaluation of remainders factors the numerator
matrix as the denominator matrix times the coefficient matrix.
The identity extends to repeated roots because both sides are
polynomials in the coefficients of the monic $Q$; polynomial
division introduces no denominators. It also includes $m=0$.
\end{proof}

\subsection*{Exact evaluation at a block boundary}
Put $N=rk$ and $c=N+r-1$. The sine formula gives
\begin{equation}\label{composition}
 P_c=A B_k.
\end{equation}
In the remainder determinant of Lemma~\ref{remainder}, retain its
first $r-1$ columns $P_N,\ldots,P_{N+r-2}$. For each
$1\le j\le r-2$, add column $P_{c-j}$ to column $P_{c+j}$.
The last $r-1$ columns then become
\[
 A B_k S_0,\ldots,A B_k S_{r-2},\qquad
 S_0=1,\quad S_j(t)=2T_j((t-a)/2)\ (j\ge1).
\]
Here $P_{c+j}+P_{c-j}=S_jP_c$, and each $S_j$ is monic of degree $j$.

Use the following monic basis of $\mathbb C[t]/(t^{r-1}A)$:
\[
 P_0,\ldots,P_{r-2},\quad A S_0,\ldots,A S_{r-2}.
\]
Its change-of-basis determinant from the monomials is one. Reduction
modulo $A$ gives an upper block with the first $r-1$ columns and zero
in the last columns. At the simple roots
$\rho_j=a+2\cos(j\pi/r)$ of $A$,
\[
 P_{N+h}(\rho_j)=(-1)^{kj}P_h(\rho_j)\qquad(0\le h\le r-2).
\]
Thus the upper block determinant is
$\prod_{j=1}^{r-1}(-1)^{kj}=\epsilon_r^k$.
The lower right block represents multiplication by $B_k$ in
$\mathbb C[t]/(t^{r-1})$. In its monomial basis this matrix is
triangular, with diagonal $B_k(0)$, so its determinant is
$B_k(0)^{r-1}$. This proves the formula for $H_N$ in
\eqref{weighted1}. No coprimality of $A$ and $t$ was assumed.

\subsection*{The adjacent determinant and the paired sum}
Temporarily choose a real $a$ for which $B_k(0)\ne0$.
This excludes only finitely many values. Equation~\eqref{composition}
gives
\[
 L(t^j B_k)=\int t^{j+r-1}P_{N+r-1}(t)\,d\mu_a(t)=0\qquad(0\le j<N).
\]
Consequently $q_N=B_k$ is the unique monic polynomial of degree $N$
orthogonal to lower degrees for $L$, since $H_N\ne0$.
The elementary identity $U_k(x)^2=\sum_{h=0}^k U_{2h}(x)$ yields
\begin{equation}\label{squared}
 A B_k^2=\sum_{h=0}^k P_{(2h+1)r-1}.
\end{equation}
For $j=0,1,2$, every $h\ge1$ term is orthogonal to $t^{r-1+j}$,
because $r\ge2$. Therefore
\begin{equation}\label{three}
 L(q_N^2)=1,\qquad L(tq_N^2)=ra,\qquad
 L(t^2q_N^2)=r+\binom{r+1}{2}a^2.
\end{equation}
For completeness, the last two values are the coefficients of
$P_{r-1}$ in $t^r$ and $t^{r+1}$. The recurrence
$tP_j=P_{j+1}+aP_j+P_{j-1}$, with $P_{-1}=0$, gives respectively
$ra$ and $\binom{r+1}{2}a^2+r$: in the latter coefficient there
are two level steps or one down step among the otherwise up steps.

The Gram determinant identity
$L(q_N^2)=H_{N+1}/H_N$ now proves $H_{N+1}=H_N$.
In particular $q_{N+1}$ also exists. Suppose first that $N\ge1$.
Let $G_N=(L(t^{i+j}))_{0\le i,j<N}$ and write
\[
 tq_N=q_{N+1}+ra\,q_N+v,\qquad \deg v<N.
\]
For $j<N-1$, $L(t^jv)=0$, whereas $L(t^{N-1}v)=1$.
Thus the coefficient vector of $v$ is $G_N^{-1}e_{N-1}$, and
\[
 L(v^2)=(G_N^{-1})_{N-1,N-1}=\frac{H_{N-1}}{H_N}.
\]
The three summands in the displayed decomposition are mutually
orthogonal. Taking squared norms and using \eqref{three} gives
\[
 r+\binom{r+1}{2}a^2
 =\frac{H_{N+2}}{H_{N+1}}+r^2a^2+\frac{H_{N-1}}{H_N},
\]
which is \eqref{weighted2}. This argument never divides by
$H_{N-1}$ or $H_{N+2}$, which are allowed to vanish.
For $N=0$, the same identity is directly
$H_2=L(t^2)-L(t)^2=r-\binom r2a^2$.
Finally, the asserted identities are polynomial identities in $a$,
so they extend from the generic real values to every complex value,
including $B_k(0)=0$. This proves Theorem~\ref{weighted}.

\subsection*{Specialization to the conjecture}
Set $a=1$, $k=n$, and $3\mid r$. Since
$T_r(-1/2)=\cos(2\pi r/3)=1$, we have $B_n(0)=U_n(1)=n+1$.
Theorem~\ref{weighted} now gives \eqref{main1}--\eqref{main2}.
For the explicit sign discrepancy, $r=9,n=1$ gives
$\epsilon_9=1$ and $H_9=H_{10}=2^8=256$.

\section*{Task 2: Sharp degrees in the remaining residue classes}
Determine the degrees and leading coefficients of the polynomial
blocks $H_{rn+s}^{(r)}$ for $2\le s<r$, including systematic zeros.

\section*{Task 3: Separate the paired determinants}
Find closed formulas for $H_{rn-1}^{(r)}$ and $H_{rn+2}^{(r)}$
individually, refining their simple sum.

\section*{Task 4: Other weights with exceptional boundaries}
Classify the real parameters $a$ for which $B_k(0)$ vanishes for
infinitely many $k$, and determine the neighboring nonzero determinants.

\section*{Task 5: A path interpretation of the exact magnitudes}
Seek a cancellation argument on families of weighted Motzkin paths
that explains the factor $(n+1)^{r-1}$ and the sign $\epsilon_r^n$.
\end{document}
